Un sifflet de train est perçu comme ayant une fréquence de \qty{\v{App}}{Hz} lorsqu’il s’approche d’un observateur immobile et de \qty{\v{Elo}}{Hz} lorsqu’il s’en éloigne. La vitesse du son est de \qty{343}{m/s}.
\begin{enumerate}[label=\alph*)]
    \item Déterminez le module de la vitesse du train, ainsi que la fréquence du son émis par le sifflet.
    \item Si le train s’immobilise, quelle est la longueur d’onde du son émis par le sifflet ?
\end{enumerate}